\chapter{Effective Field Theories Without Poincare-Invariant UV Completion}
\label{ch:efts-without-poincare-uv-completion}
\ConventionsEuclidean

\section{Scaling Limits and Finite-Cutoff Effective Theories}

This chapter separates two constructions that are often conflated.  A
lattice or many-body system may have a critical scaling limit which is a
continuum QFT, or it may have only a finite-cutoff long-distance effective
description.  The first case can produce a UV-complete relativistic theory.
The second case is an effective theory whose validity is tied to the
microscopic cutoff and to a specified range of observables.

\begin{definition}[Critical scaling limit]
\label{def:critical-scaling-limit}
Let \(a>0\) be a lattice spacing, let
\(\mu_{a,\kappa}\) be a finite- or infinite-volume Euclidean lattice
probability measure with couplings \(\kappa\), and let
\(\mathcal O_{a,i}(x)\) be lattice observables.  A critical scaling limit
consists of:
\begin{enumerate}
\item a tuning \(\kappa(a)\) for which the correlation length in lattice
      units diverges;
\item normalization factors \(Z_i(a)\) and continuum coordinates
      \(x=a n\);
\item limits of separated Euclidean correlators
      \[
        \lim_{a\downarrow0}
        \left\langle
          \prod_{j=1}^k Z_{i_j}(a)\mathcal O_{a,i_j}(n_j)
        \right\rangle_{\mu_{a,\kappa(a)}} ;
      \]
\item Euclidean covariance, locality, reflection positivity, clustering, and
      regularity estimates sufficient for OS or another stated
      reconstruction theorem.
\end{enumerate}
When these data are supplied, the limiting continuum QFT is the scaling-limit
object constructed from the lattice system.
\end{definition}

\begin{definition}[Finite-cutoff effective theory]
\label{def:finite-cutoff-effective-theory}
A finite-cutoff effective theory is a tuple
\[
  (\Lambda_{\rm UV},\mathcal A_\Lambda,S_\Lambda,\mathcal O_\Lambda,
   \mathcal R,E_{\rm max})
\]
where \(\Lambda_{\rm UV}\) is a physical ultraviolet scale,
\(\mathcal A_\Lambda\) is the algebra or field-variable space retained below
that scale, \(S_\Lambda\) is the regulated action or Hamiltonian,
\(\mathcal O_\Lambda\) is the declared operator-insertion prescription,
\(\mathcal R\) is the set of observables to be matched, and
\(E_{\rm max}\ll\Lambda_{\rm UV}\) is the energy or momentum domain in which
the error estimate is claimed.  It is a QFT technique only after the
regulator, matching observables, and error class have been stated.
\end{definition}

The two definitions have different logical content.  A critical scaling limit
aims to remove the microscopic cutoff.  A finite-cutoff effective theory keeps
the cutoff as part of its meaning, even when its formulas resemble a
relativistic continuum Lagrangian.

\begin{proposition}[Scaling limit as continuum QFT]
\label{prop:scaling-limit-continuum-qft}
Suppose the limits in Definition~\ref{def:critical-scaling-limit} exist for a
set of local lattice observables closed under the operator products needed to
generate the limiting Schwinger functions, and suppose the limiting
Schwinger functions satisfy the OS hypotheses with Euclidean covariance.
Then the reconstructed theory is a continuum local QFT.  If the limiting
Euclidean group enhances to the full Euclidean rotation group and the OS
linear-growth hypothesis needed for continuation is verified, the Lorentzian
reconstruction is a Poincare-covariant Wightman theory.
\end{proposition}

\begin{proof}
The OS reconstruction theorem applies to the limiting Schwinger functions,
not to the finite lattice variables themselves.  Reflection positivity gives a
positive Hilbert-space inner product after quotienting null vectors;
Euclidean time translations give a positive Hamiltonian; Euclidean covariance
and the required analytic-growth bound give the Lorentzian representation of
the Poincare group; and locality follows from Euclidean symmetry and the OS
permutation/reflection structure after analytic continuation.  The
microscopic lattice remains a regulator for the construction, while the
Hilbert space, fields, and correlation functions are those of the limiting
object.
\end{proof}

Thus the two-dimensional Ising scaling limit and, conjecturally, the
three-dimensional Ising critical point belong to the first category.  Away
from criticality, a gapped lattice system may still be described at long
distances by a continuum effective action, but the continuum action is then a
finite-cutoff description of selected observables, not a removed-cutoff
definition of a Poincare-invariant QFT.

\section{What Must Be Revalidated}

Let \(\mathcal T_{\rm rel}\) be a theorem whose proof in earlier volumes used
Poincare covariance, the spectral condition, local commutativity, and the
existence of asymptotic particle states.  Applying the same theorem to a
condensed-matter EFT with microscopic cutoff \(a^{-1}\) requires a replacement
proof or a matching argument.

\begin{definition}[Validation map for borrowed QFT technology]
\label{def:validation-map-borrowed-qft}
A validation map from a relativistic theorem or construction
\(\mathcal T_{\rm rel}\) to a finite-cutoff effective theory consists of:
\begin{enumerate}
\item the precise hypotheses \(H_1,\ldots,H_n\) used in the relativistic
      proof;
\item finite-cutoff replacements \(H_i^\Lambda\) or a proof that \(H_i\) is
      emergent with a stated error bound;
\item the observable class \(\mathcal R\) on which the conclusion is claimed;
\item an estimate comparing the microscopic observable and the effective
      observable in that class;
\item a statement of which parts of the relativistic conclusion are retained,
      modified, or unavailable.
\end{enumerate}
\end{definition}

The following table records typical outcomes.
\begin{center}
\small
\begin{tabular}{@{}p{0.31\textwidth}p{0.58\textwidth}@{}}
\hline
Relativistic tool & Finite-cutoff status\\
\hline
Wilsonian integration &
Available after the cutoff, fields, and blocking map are stated.\\
Power counting &
Available with the correct dynamical exponent and symmetries.\\
OS/Wightman analyticity &
Requires an independent Euclidean-to-real-time argument.\\
LSZ scattering &
Requires stable quasiparticles and an asymptotic dynamics.\\
Spin-statistics and CPT &
Do not transfer without Lorentz/locality hypotheses.\\
Wedge modular flow &
Requires a relativistic wedge or a replacement modular theorem.\\
Anomaly matching &
Transfers only for the specified symmetry and regulator response.\\
\hline
\end{tabular}
\end{center}
The entries are not restrictions on using effective field theory.  They are
the data needed to know what the calculation proves.

\section{Examples}

\subsection{Fermi Liquid}

A Fermi liquid is organized around a codimension-one Fermi surface
\(\Sigma_F\) in momentum space, not around the Lorentz-invariant point
\(p_\mu p^\mu=m^2\).  A local patch with outward normal \(\hat n\) has
linearized dispersion
\[
  \epsilon(k_F\hat n+\ell\hat n+q_\parallel)
  =
  v_F\ell+O(\ell^2,q_\parallel^2),
\]
where \(q_\parallel\) is tangent to \(\Sigma_F\).  The scaling that keeps the
kinetic term fixed scales \(\ell\) and frequency differently from tangential
momenta.  Hence relativistic spin-statistics, crossing, and LSZ are not the
primitive logic.  The primitive data are the microscopic Hamiltonian, the
Fermi surface, quasiparticle lifetime estimates, Landau parameters, and the
range of frequencies and momenta for which quasiparticles are sharp.

\subsection{Lifshitz Fixed Points}

A Lifshitz scalar with dynamical exponent \(z\) has a quadratic action of the
form
\[
  S=\frac12\int\dd\tau\,\dd^d x\,
  \left[
    (\partial_\tau\phi)^2+\kappa^2(\Delta^{z/2}\phi)^2
  \right]
\]
when \(z\) is even, with the obvious Fourier-symbol replacement for general
integer \(z\).  The scaling is
\[
  x\mapsto b x,\qquad \tau\mapsto b^z\tau .
\]
Power counting is meaningful, but it is not relativistic power counting.
If a Lorentz-invariant relativistic action appears in the infrared, the
Lorentz-violating operators must be shown to be irrelevant or tuned small in
the stated scaling regime.

\subsection{Quantum Hall Effective Theory}

The integer or fractional quantum Hall response at long distances may be
encoded by a Chern--Simons effective action
\[
  S_{\rm CS}[a,A]
  =
  \frac{K_{IJ}}{4\pi}\int a^I\dd a^J
  +\frac{q_I}{2\pi}\int A\dd a^I+\cdots .
\]
This is a topological response theory for a gapped many-body phase with a
microscopic nonrelativistic Hamiltonian and background electromagnetic field.
Its gauge invariance, boundary anomaly, and line-operator data are real
low-energy statements, but they do not imply a local Poincare-invariant UV
completion.  The correct validation map matches ground-state response,
edge anomaly, quasiparticle braiding, and finite-size gaps, not an
asymptotic-particle S-matrix.

\section{Controlled Use of Emergent Relativistic Language}

\begin{controlledapproximation}[Emergent relativistic EFT]
\label{ca:emergent-relativistic-eft}
Let a microscopic system with cutoff \(\Lambda_{\rm UV}\) have low-energy
observables \(\mathcal R\).  A relativistic EFT with action \(S_{\rm rel}\)
is a controlled emergent description on \(\mathcal R\) if there exists a
matching map \(M_\Lambda\) and constants \(C,\Delta>0\) such that for every
observable \(O\in\mathcal R\) with characteristic scale \(E\le E_{\rm max}\),
\[
  \left|
    \langle O\rangle_{\rm micro}
    -
    \langle M_\Lambda O\rangle_{S_{\rm rel}}
  \right|
  \le
  C_O\left(\frac{E}{\Lambda_{\rm UV}}\right)^{\Delta}
\]
after the stated renormalizations.  The exponent \(\Delta\) is the smallest
scaling dimension gap between retained operators and omitted symmetry-allowed
Lorentz-violating operators in the chosen RG chart.
\end{controlledapproximation}

This formulation is deliberately observable-based.  It does not require that
the microscopic system possess a Lorentzian spacetime symmetry.  It requires
only that the selected observables are approximated by a relativistic
continuum theory with a quantitative error bound.  In contrast, a critical
scaling limit satisfying Proposition~\ref{prop:scaling-limit-continuum-qft}
constructs the continuum QFT as its limiting object.

\begin{openproblem}[Validated condensed-matter EFT atlas]
\label{op:validated-condensed-matter-eft-atlas}
For each major class of condensed-matter effective theory, construct an
atlas of validation maps in the sense of
Definition~\ref{def:validation-map-borrowed-qft}: Fermi liquids, non-Fermi
liquids, superconductors, deconfined critical points, fractional quantum Hall
states, spin liquids, fracton phases, and systems with emergent gauge fields.
The goal is not to force these systems into Poincare-invariant axioms, but to
state exactly which QFT techniques are justified for which observables.
\end{openproblem}
